\begin{explanationbox}
    \textbf{\textcolor{CExplanation}{把两条直角边一起转过去}}
\begin{center}
\begin{tikzpicture}[x=0.65cm,y=0.65cm,>=stealth,every node/.style={font=\small},line cap=round]

\coordinate (B) at (0,0); \coordinate (P) at (3,2); \coordinate (Q) at (-2,3);
\fill[blue!10] (B)--(3,0)--(P)--cycle;
\fill[teal!10] (B)--(0,3)--(Q)--cycle;
\draw[->,gray] (-3.7,0)--(4.5,0);
\draw[->,gray] (0,-0.7)--(0,4.2);
\draw[very thick,blue!75!black] (B)--(P);
\draw[very thick,teal!75!black] (B)--(Q);
\draw[thick,orange!85!black,->] (B)--(3,0) node[midway,below] {$u$};
\draw[thick,purple,->] (3,0)--(P) node[midway,right] {$v$};
\draw[thick,orange!85!black,->] (B)--(0,3) node[pos=.65,right] {$u$};
\draw[thick,purple,->] (0,3)--(Q) node[midway,above] {$-v$};
\draw (2.7,0)--(2.7,.3)--(3,.3);
\draw (0,2.7)--(-.3,2.7)--(-.3,3);
\node[below left] at (B) {$B$};
\node[above right] at (P) {$P:(u,v)$};
\node[above left] at (Q) {$Q:(-v,u)$};
\node[anchor=west] at (6.2,2.8) {逆时针旋转};
\node[anchor=west] at (6.2,1.8) {水平 $u$ 变竖直 $u$};
\node[anchor=west] at (6.2,.8) {竖直 $v$ 变水平 $-v$};
\node[anchor=west] at (-3.5,-1.2) {图示 $u>0,v>0$；负值表示沿相反方向。};

\end{tikzpicture}
\end{center}


    看同色的边：水平变化 $u$ 转为竖直变化 $u$，竖直变化 $v$ 转为水平变化 $-v$。
    因此逆时针是 $(u,v)\mapsto(-v,u)$；反向转则是 $(u,v)\mapsto(v,-u)$。
    图示取 $u,v>0$；若其中一个为负，只需沿相反方向走，公式不变。

    \medskip
    \textbf{\textcolor{CExplanation}{从位移回到点坐标}}
    绕 $B(a,b)$ 旋转，先求 $P-B=(u,v)$，再旋转这个位移，最后加回 $B$。
    在证明方法二的实坐标图中，可以直接核对这三步。

    \medskip
    \textbf{\textcolor{CExplanation}{先用交换关系消去一个未知量}}
    若 $Q$ 在 $x=a+d$ 上，它相对 $B$ 的横向变化已固定为 $d$，故 $|y_1-b|=|d|$。
    后面的椭圆例题把这两个相等的横距、竖距直接标在图上。
\end{explanationbox}
